// 323111047

\section{Introducción}

Planteamiento del problema. El reto técnico de esta actividad consiste en diseñar una aplicación de consola en C que simule un reproductor de audio mediante una lista ligada circular. El problema principal radica en la gestión de punteros para asegurar que el último nodo de la lista apunte siempre al primero, eliminando la existencia de un puntero nulo al final de la estructura. Se deben implementar funciones para la inserción, búsqueda y eliminación de canciones, además de permitir la navegación continua a través de la lista, lo cual exige un control preciso de la memoria dinámica para mantener la integridad de los datos durante las modificaciones de la lista de reproducción.

\vspace{0.4cm}

Motivación. La importancia de resolver este problema se basa en la necesidad de comprender estructuras de datos que permiten flujos de información cíclicos, los cuales son fundamentales en aplicaciones multimedia y sistemas operativos. En el desarrollo de software, las listas circulares son ideales para procesos que requieren repetición continua sin interrupciones, como las listas de reproducción de audio o el despacho de tareas en un procesador. Estudiar este modelo permite al estudiante optimizar el acceso a los datos y desarrollar una lógica de programación más avanzada para el manejo de estructuras ligadas complejas que no poseen un final lineal.

\vspace{0.4cm}

Objetivos. Se busca que el alumno demuestre dominio en las operaciones de inserción y eliminación de nodos sin romper la circularidad de la estructura, así como en la búsqueda eficiente de elementos por nombre. Finalmente, se pretende validar que el programa sea capaz de avanzar entre canciones de forma infinita y simular la reproducción de las mismas, proporcionando una base sólida para discutir el cumplimiento de los requerimientos técnicos en la sección de resultados.